En somme, ce projet professionnel n'est pas une simple accumulation de diplômes, 
mais la structuration d'une trajectoire cohérente. 

En articulant la rigueur de la physique fondamentale, 
la méthode académique du Magistère de Saclay 
et l'ouverture qu'offre une formation généraliste, 
je me donne les moyens de passer de la théorie à l'application concrète. 

L'ingénierie que je vise est avant tout un métier de traduction : 
celle qui permet de convertir un problème complexe en une solution modélisable et viable. 

C'est cette perspective, alliant précision scientifique et pragmatisme, 
qui guidera mes choix dans les prochaines années, 
que ce soit dans le paysage industriel français ou à l'international.